\newpage
\section{Ergebnisse und Diskussion}\label{diskussion}

Befund und Deutung bleiben in diesem Kapitel getrennt. Die Ergebnisse der neun Ablationsläufe, der Segmentierung und der Verbuschungsberechnung werden zunächst berichtet, ihre Einordnung folgt geschlossen im Diskussionsteil.

\subsection{Ergebnisse der Ablationsstudie}\label{ergebnis_ablation}

Getrennt nach den beiden Testmengen, die in Abschnitt~\ref{ablation} erläutert werden, halten die Tabellen~\ref{tab:ablation_a} und~\ref{tab:ablation_b} die Kennzahlen aller neun Läufe fest. Aussagekräftig ist jeweils nur der Vergleich innerhalb einer Tabelle, wobei run07 als einziger Lauf auf beiden Testmengen ausgewertet vorliegt und die beiden Teile miteinander verbindet. Von den acht Eingriffen, die gegenüber ihrem Vergleichslauf genau einen Faktor ändern, wirkt einer.

\begin{table}[H]
    \centering
    \caption[Ergebnisse der Läufe run01 bis run04]{Ergebnisse der Läufe run01
    bis run04 auf der Testmenge mit 487 Kacheln. Die mIoU von run04 mittelt
    über sechs statt acht Klassen. Eigene Darstellung.}
    \label{tab:ablation_a}
    \begin{tabular}{lrrrrr}
    \hline
    Lauf & Klassen & OA & mIoU & IoU bush & Val.-mIoU (Ep.) \\
    \hline
    run01 & 8 & 0,671 & 0,248 & 0,394 & 0,356 (53) \\
    run02 & 8 & 0,674 & 0,242 & 0,339 & 0,349 (7)  \\
    run03 & 8 & 0,650 & 0,229 & 0,329 & 0,346 (6)  \\
    run04 & 6 & 0,662 & 0,344 & 0,393 & 0,447 (58) \\
    \hline
    \end{tabular}
\end{table}

Der Nahinfrarotkanal bringt in der Gesamtbilanz keinen Gewinn, denn mittlere IoU und Gesamtgenauigkeit bleiben unverändert. Zwischen den Klassen verschiebt sich lediglich die Güte, denn die IoU von \texttt{graminoid} steigt von 0,292 auf 0,375, wohingegen \texttt{heath} von 0,262 auf 0,177 zurückfällt. Aufschlussreicher als die Kennzahlen selbst ist der Zeitpunkt des besten Modells, das statt in Epoche 53 bereits in Epoche 7 erreicht wird, denn ein Kanal ohne zusätzlichen Informationsgehalt führt das Netz früher in die Überanpassung. Naheliegend ist, dass sich die grünen Vegetationsklassen auch im Nahinfrarot zu ähnlich verhalten und gerade die Trennung ausbleibt, auf die es ankommt.

Noch deutlicher fällt das Urteil über den Focal Loss aus, der klassenweise durchweg schlechtere Werte liefert und auch die Gesamtgenauigkeit senkt. Am härtesten trifft es die ohnehin schwachen Klassen, denn die IoU von \texttt{soil} bricht von 0,035 auf 0,001 ein und \texttt{other} verschwindet ganz. Zusammen mit den bereits kräftigen Klassengewichten überbetont der Focal Loss offenbar die seltenen Klassen, die sich in diesem Datensatz grundsätzlich nicht trennen lassen und verschlechtert im selben Zug die Klassen, die schon vorher gut trennbar waren.

Der Wechsel auf sechs Klassen hebt die mittlere IoU von 0,242 auf 0,344, doch der größere Teil dieses Sprungs ist mechanisch bedingt. Bildet man die Vorhersagen von run02 nachträglich auf dasselbe Schema ab, indem \texttt{sand} und \texttt{soil} zu \texttt{ground} zusammengefasst und die Klasse \texttt{other} entfernt wird, so ergibt sich bereits eine mittlere IoU von 0,336. Auf das Training mit weniger Klassen entfällt somit nur der kleinere Rest. Aufschlussreich ist trotzdem die Klasse \texttt{ground}, die durch das bloße Zusammenfassen von einer IoU um 0,04 auf 0,162 springt, denn Sand und offener Boden wurden gegeneinander verwechselt und nicht grundsätzlich übersehen. Für die Fragestellung dieser Arbeit günstig ist außerdem, wie sich die Güte zwischen den Vegetationsklassen verschiebt, denn die IoU von \texttt{bush} steigt von 0,339 auf 0,393 und die von \texttt{heath} von 0,177 auf 0,255, zulasten von \texttt{graminoid}, das von 0,375 auf 0,258 zurückfällt. Die Kernverwechslung löst die neue Taxonomie nicht, wohl aber erlaubt sie ein stabileres Training, wie das späte Optimum in Epoche 58 statt in Epoche 6 oder 7 zeigt.

Von den beiden lokalen Höhenmerkmalen war zu erwarten, dass sie Büsche und gleich hohe Heide voneinander lösen, doch die mittlere IoU steigt nur von 0,344 auf 0,348. Nach wie vor sagt das Netz gut vier von zehn Buschpunkten als \texttt{graminoid} voraus. Verantwortlich ist die Auflösung der Merkmale, denn \texttt{z\_range} und \texttt{z\_rel} beruhen auf Rasterzellen von 0,5 beziehungsweise 2,5 Metern, sodass alle Punkte einer Zelle denselben Bezugswert erhalten und ein Grasbüschel neben einem Busch dessen Höhenspanne erbt. Genau an dieser Stelle sitzt der Fehler, den die Merkmale beheben sollten.

\begin{table}[H]
    \centering
    \caption[Ergebnisse der Läufe run05 bis run09]{Ergebnisse der Läufe run05
    bis run09 auf der Testmenge mit 468 Kacheln. Die mIoU von run06 steht in
    Klammern, weil sie über zwei statt sechs Klassen mittelt.
    Eigene Darstellung.}
    \label{tab:ablation_b}
    \begin{tabular}{lrrrrrrr}
    \hline
    Lauf & Klassen & mIoU & IoU bush & IoU Gehölz & Precision & Recall & F1 \\
    \hline
    run05 & 6 & 0,357 & 0,413 & 0,724 & 0,934 & 0,763 & 0,840 \\
    run06 & 2 & (0,725) & -- & 0,713 & 0,966 & 0,731 & 0,833 \\
    run07 & 6 & \textbf{0,393} & \textbf{0,500} & \textbf{0,761} & 0,930 & 0,808 & \textbf{0,864} \\
    run08 & 6 & 0,361 & 0,452 & 0,730 & 0,892 & 0,801 & 0,844 \\
    run09 & 6 & 0,347 & 0,359 & 0,701 & 0,942 & 0,733 & 0,824 \\
    \hline
    \end{tabular}
\end{table}

Mit run06 wird das Klassenschema ein zweites Mal vereinfacht, diesmal bis auf zwei Klassen, sodass das Netz von Grund auf nur noch Gehölz gegen alles Übrige lernt. Dahinter stand die Überlegung, dass ein Teil der Kapazität auf die schwierige Trennung von Heide und Gras entfällt, die für den Verbuschungsgrad keine Rolle spielt. Gemessen wird es deshalb am einfacheren Weg, die sechs Klassen von run05 im Nachhinein zusammenzufassen (Abschnitt~\ref{evaluation}). Das gelingt nicht, denn das binäre Modell bleibt bei Gehölz-IoU und F1-Wert knapp hinter dem zusammengefassten Sechsklassenmodell zurück und verschiebt lediglich das Verhältnis von Precision und Recall, indem es die Klasse Gehölz seltener vergibt und dabei seltener irrt. Die Vermutung, das Netz verschwende Kapazität, bestätigt sich somit nicht. Die Grenze zwischen Gehölz und niedriger Vegetation ist in den Daten selbst schwach ausgeprägt und wird nicht erst durch das Klassenschema verdeckt.

Wirkung zeigt einzig die Vergrößerung der Nachbarschaftsradien, die aus der in Abschnitt~\ref{ablation} beschriebenen Messung der Punktdichte hervorging. Auf der Testmenge mit 468 Kacheln steigen alle für die Fragestellung wichtigen Kennzahlen deutlich, die IoU von \texttt{bush} von 0,413 auf 0,500 und die Gehölz-IoU von 0,724 auf 0,761. Auf der ursprünglichen Testmenge mit 487 Kacheln sinkt die mittlere IoU im selben Vergleich dagegen von 0,348 auf 0,336, weil \texttt{graminoid} und \texttt{ground} verlieren. In dieser Gegenläufigkeit zeigt sich ein Skalenkonflikt, denn ein Busch ist ein großes Objekt und braucht Kontext, wohingegen die Unterscheidung von Gras und Zwergstrauchheide an feiner Textur hängt und mit größeren Nachbarschaften verwischt wird. Für diese Arbeit ist der Tausch günstig, weil sie vor allem eine verlässliche Gehölzerkennung braucht.

Die zweite mögliche Antwort auf die geringe Punktdichte, die verdoppelte Punktzahl je Kachel, verschlechtert das Ergebnis auf identischer Testmenge durchweg, wobei sich einzig \texttt{ground} verbessert. Auch Multi-Scale Grouping bringt keinen Gewinn und bleibt der schwächste Lauf der Reihe, wobei einzig \texttt{graminoid} den zweitbesten Wert aller Läufe erreicht. Sein zusätzlicher feiner Zweig arbeitet mit einem Radius von 0,125 Metern und findet nach der gemessenen Punktdichte selbst bei 8.192 Punkten je Kachel nur ein bis zwei Nachbarn, sodass er keine zusätzliche Information über einen Punkt liefert.

Über alle Läufe hinweg blieb die Verwechslung zwischen Zwergstrauchheide, Gras und Busch bestehen, obwohl mit dem Nahinfrarot, dem Focal Loss, dem reduzierten Klassenschema, den expliziten Höhenmerkmalen und dem binären Modell fünf voneinander unabhängige Ansätze an ihr angesetzt haben. Dass am Ende ein Parameter der Nachbarschaftssuche wirkt und kein zusätzliches Merkmal, spricht dafür, dass dem Netz nicht Informationen fehlte, sondern eine zur Punktdichte passende räumliche Bezugsgröße. Bemerkenswert ist daneben die Klasse \texttt{tree}, deren IoU in allen acht Mehrklassenläufen zwischen 0,937 und 0,963 liegt. Strukturell eindeutige Objekte werden zuverlässig erkannt und die Schwäche betrifft ausschließlich die niedrige Vegetation.

Für alle Vergleiche in beiden Tabellen gilt allerdings eine Einschränkung, die aus dem Zuschnitt des Validierungsteils folgt. Die besten Validierungswerte von run07, run08 und run09 liegen bei 0,4684, 0,4685 und 0,4656 und somit praktisch gleichauf. Auf dem Testteil streuen dieselben Modelle dagegen zwischen 0,347 und 0,393. Der Validierungsteil besteht aus drei räumlich getrennten Blöcken, von denen ein einziger 378 der 455 Kacheln stellt, sodass das Auswahlkriterium im Wesentlichen an einer Teilfläche hängt. Welche der drei Varianten als bestes Modell gespeichert wird, entscheidet sich deshalb zu einem erheblichen Teil zufällig und Unterschiede von wenigen Tausendsteln sind nicht als Effekt zu deuten.

Als finales Modell dient deshalb run07. Es erreicht auf der Testmenge mit 468 Kacheln zugleich die beste mittlere IoU und die beste Gehölzerkennung. Seine Radien gehen zudem auf die gemessene Punktdichte der Kacheln zurück. Abschnitt~\ref{evaluation_segmentierung} betrachtet seine Ergebnisse im Einzelnen.

\subsection{Güte der semantischen Segmentierung}\label{evaluation_segmentierung}

Ausgewertet wird run07 auf dem Testteil der annotierten Kacheln, der während
des gesamten Trainings unberührt geblieben ist. Dort erreicht das Modell eine
Gesamtgenauigkeit von 0,650 und eine mittlere IoU von 0,393. Hinter diesen
beiden Werten steht eine große Spannweite, die Tabelle~\ref{tab:metriken}
klassenweise aufschlüsselt. Bäume erkennt das Modell nahezu fehlerfrei,
\texttt{bush} erreicht ein mittleres Niveau und \texttt{ground} sowie
\texttt{deadwood} bleiben schwach.

\begin{table}[H]
    \centering
    \caption[Klassenweise Kennzahlen des Modells run07]{Klassenweise Kennzahlen
    des Modells run07 auf dem Testteil. Der Anteil gibt die wahren Punkte der
    Klasse an. Eigene Darstellung.}
    \label{tab:metriken}
    \begin{tabular}{lrrrrr}
    \hline
    Klasse & IoU & Precision & Recall & F1 & Anteil \\
    \hline
    bush      & 0,500 & 0,788 & 0,577 & 0,666 & 25,9 \% \\
    deadwood  & 0,196 & 0,239 & 0,523 & 0,328 & 1,2 \%  \\
    graminoid & 0,296 & 0,583 & 0,376 & 0,457 & 28,5 \% \\
    heath     & 0,322 & 0,358 & 0,762 & 0,487 & 14,9 \% \\
    ground    & 0,103 & 0,518 & 0,114 & 0,187 & 1,3 \%  \\
    tree      & 0,938 & 0,971 & 0,965 & 0,968 & 28,3 \% \\
    \hline
    Gehölz    & 0,761 & 0,930 & 0,808 & 0,864 & 54,2 \% \\
    \hline
    \end{tabular}
\end{table}

Maßgeblich für die Fragestellung ist die letzte Zeile dieser Tabelle. Auf
der Ebene von Gehölz gegen Rest findet das Modell vier von fünf tatsächlich
vorhandenen Gehölzpunkten (Recall). Von den als Gehölz ausgewiesenen Punkten gehören über neun von zehn wirklich dazu (Precision). Die mittlere IoU fällt vor allem deshalb so niedrig aus, weil ein erheblicher Teil der Fehler zwischen \texttt{graminoid} und \texttt{heath} entsteht.

Wo genau die Fehler auftreten, zeigt die Konfusionsmatrix in Tabelle~\ref{tab:konfusion}. Mehr als die Hälfte der Graspunkte weist das Modell der Zwergstrauchheide zu, die falsch zugeordneten Buschpunkte verteilen sich auf \texttt{graminoid} und \texttt{heath}, und offener Boden geht zu mehr als zwei Dritteln an \texttt{heath} verloren. Sämtliche dieser Verwechslungen betreffen niedrige Bestände, die sich in Höhe und Farbe überschneiden. Bäume dagegen ordnet das Modell zu 96,5 Prozent richtig zu. Die wenigen Fehler betreffen fast nur die Klasse \texttt{bush}.

\begin{table}[H]
    \centering
    \caption[Konfusionsmatrix des Modells run07]{Konfusionsmatrix des Modells
    run07 in Prozent der wahren Klasse. Die Hauptdiagonale entspricht der
    Accuracy je Klasse. Eigene Darstellung.}
    \label{tab:konfusion}
    \begin{tabular}{lrrrrrr}
    \hline
    wahr $\backslash$ vorhergesagt & bush & deadw. & gramin. & heath & ground & tree \\
    \hline
    bush      & \textbf{57,7} & 3,0  & 20,9 & 15,4 & 0,0  & 3,0 \\
    deadwood  & 16,1 & \textbf{52,3} & 12,5 & 14,4 & 1,9  & 2,9 \\
    graminoid & 7,8  & 0,5  & \textbf{37,6} & 54,0 & 0,1  & 0,0 \\
    heath     & 5,8  & 5,1  & 12,5 & \textbf{76,2} & 0,4  & 0,0 \\
    ground    & 1,4  & 0,5  & 18,6 & 68,0 & \textbf{11,4} & 0,0 \\
    tree      & 2,6  & 0,9  & 0,0  & 0,0  & 0,0  & \textbf{96,5} \\
    \hline
    \end{tabular}
\end{table}

Die niedrigsten IoU-Werte erreichen mit \texttt{deadwood} und \texttt{ground}  zugleich die beiden seltensten Klassen, denn sie stellen zusammen nur gut zwei Prozent der Testpunkte. Trotz der Gewichtung im Verlust (Abschnitt~\ref{verlustfunktion}) lernt das Netz sie kaum. Da die mittlere IoU alle Klassen gleich gewichtet, ziehen sie den Mittelwert spürbar nach unten. Für die Zielgröße
bleibt das folgenlos, weil weder Totholz noch offener Boden in den Verbuschungsgrad eingehen.

\subsection{Anwendung auf die Befliegungsgebiete}\label{ergebnis_inferenz}

Auf dieser Grundlage wurde run07 auf vier Befliegungen im Authausener Wald
angewendet, nämlich die 110-kV-Trasse mit dem angrenzenden PSA-Platz, die
HeideWildnis-Punktwolke sowie die Schneise I und die Schneisen IV und V. Mit 67,4
Millionen Punkten ist die 110-kV-Trasse das mit Abstand größte dieser Gebiete,
die übrigen drei liegen zwischen 8,3 und 12,6 Millionen. Klassifiziert wurden
auf diese Weise 65,9 Hektar belegter Fläche, für die zuvor keine Annotation
vorlag.

\begin{figure}[H]
    \centering
    \includegraphics[width=0.75\textwidth]{vegetationsklassen_heidewildnis}
    \caption[Vorhergesagte Vegetationsklassen im Gebiet HeideWildnis]{Vom
    Modell run07 vorhergesagte Vegetationsklassen im Gebiet HeideWildnis,
    weiße Umrisse: Polygone der Biotopkartierung. Eigene Darstellung,
    Hintergrund: Google Satellite.}
    \label{fig:vegetationsklassen_heidewildnis}
\end{figure}

Wie eine solche Vorhersage aussieht, zeigt
Abbildung~\ref{fig:vegetationsklassen_heidewildnis} am Beispiel der
HeideWildnis-Punktwolke. Die Offenfläche erscheint als zusammenhängender Bestand aus
Zwergstrauchheide und Gras, durchsetzt von einzelnen Buschgruppen, und wird
ringsum von einem geschlossenen Saum der Klasse \texttt{tree} eingefasst.

Über ein ganzes Gebiet hinweg lässt sich das in Klassenanteilen fassen, hier
stellvertretend für die 110-kV-Trasse als größtes der vier Gebiete. Knapp die
Hälfte aller Punkte entfällt auf die Klasse \texttt{tree}, was der bewaldeten
Umgebung der Trasse entspricht, gut ein Zehntel auf \texttt{bush} und der
übrige Anteil überwiegend auf \texttt{graminoid} und \texttt{heath}.
Aussagekräftiger als die Verteilung selbst ist die je Punkt mitgeführte
Sicherheit, denn sie ordnet sich in derselben Reihenfolge wie die Güte auf dem
Testteil. Bei Bäumen liegt sie im Mittel bei 0,98, bei \texttt{deadwood} mit
0,58 am tiefsten. Unsicher ist sich das Modell genau dort, wo es nachweislich
häufig falsch liegt.

Ohne Vorhersage bleiben 2,25 Prozent der Punkte, die so dünn belegte Bereiche
besetzen, dass kein sie überdeckendes Fenster die Mindestzahl von hundert
Punkten erreicht. Für die Auswertung fällt das kaum ins Gewicht, weil ein
Kartierungspolygon ohnehin nur dann einbezogen wird, wenn mindestens 80
Prozent seiner Fläche mit Vorhersagen belegt sind
(Abschnitt~\ref{validierung}).

Wie belastbar diese Vorhersagen sind, hängt daran, wo die annotierten
Teilflächen liegen. Acht der elf Authausener Teilflächen liegen in der
Schneise I, die übrigen drei in der 110-kV-Trasse. In der
HeideWildnis-Punktwolke und in den Schneisen IV und V hat das Modell folglich
nie trainiert, sodass jede dortige Vorhersage, auch die in
Abbildung~\ref{fig:vegetationsklassen_heidewildnis} gezeigte, eine Übertragung
auf unbekanntes Gelände ist und für die Frage der Übertragbarkeit am meisten
aussagt. Klassifiziert wird trotzdem jedes Gebiet lückenlos, denn der
Verbuschungsgrad soll gerade für die Flächen außerhalb der annotierten
Ausschnitte entstehen. Diese Ausschnitte bedecken rund fünf Prozent der
klassifizierten Fläche, in der kleinen Schneise I allerdings über die Hälfte.

\subsection{Berechneter Verbuschungsgrad im Vergleich zur Kartierung}\label{ergebnis_quantifizierung}

Als Referenz gehen 43 Offenlandpolygone der Kartierung von 2014 in die Auswertung ein, die zusammen 24,9 Hektar bedecken und die in Abschnitt~\ref{validierung} genannten Bedingungen erfüllen. Dem Lebensraumtyp 4030 der trockenen europäischen Heiden gehören 14 von ihnen mit zusammen 9,9 Hektar an. Sieben Polygone tragen im Bemerkungsfeld einen Hinweis auf eine kurz vor der Kartierung erfolgte Pflegemaßnahme. Die Kennzahlen der folgenden Abschnitte beruhen auf den 14 Heidepolygonen mit zusammen 9,9 Hektar, der Vergleich der Höhengrenzen nur auf den sieben Flächen ohne Pflegevermerk mit 5,8 Hektar. Gemessen an den 64,5 Hektar, für die das Verfahren einen Wert liefert, sind das 15 beziehungsweise 9 Prozent.

\begin{table}[H]
    \centering
    \caption[Vergleich der Verbuschungsdefinitionen]{Übereinstimmung des
    berechneten mit dem kartierten Verbuschungsgrad für sechs Definitionen,
    gerechnet über die sieben Heidepolygone ohne Pflegevermerk. RMSE und
    Versatz in Prozentpunkten. Eigene Darstellung.}
    \label{tab:hoehengrenze}
    \begin{tabular}{lrrrrr}
    \hline
    Definition & RMSE & Versatz & $R^2$ & Pearson & Spearman \\
    \hline
    nur \texttt{bush}          & 15,5 & $-6,9$ & 0,07 & 0,53 & 0,45 \\
    Gehölz ohne Grenze         & 16,9 & $+9,2$ & $-0,10$ & 0,72 & 0,04 \\
    Gehölz unter 2 m           & 16,4 & $-7,6$ & $-0,03$ & 0,46 & 0,45 \\
    Gehölz unter 3 m           & 12,4 & $-4,6$ & 0,41 & 0,72 & 0,38 \\
    \textbf{Gehölz unter 5 m}  & \textbf{9,5} & \textbf{$+1,8$} & \textbf{0,65} & \textbf{0,86} & 0,23 \\
    Gehölz unter 8 m           & 11,0 & $+3,8$ & 0,54 & 0,83 & 0,34 \\
    \hline
    \end{tabular}
\end{table}

Welche Höhengrenze die beste Übereinstimmung liefert, zeigt Tabelle~\ref{tab:hoehengrenze}. Die Grenze von fünf Metern erreicht mit 9,5 Prozentpunkten den kleinsten mittleren Fehler und mit $+1,8$ Prozentpunkten zugleich den geringsten Versatz. Ohne jede Begrenzung überschätzt das Verfahren die Verbuschung deutlich, weil die überhängenden Kronen des angrenzenden Waldes mit in die Berechnung eingehen. Im Polygon 13001\_2623 etwa werden 71,7 Prozent der Fläche als \texttt{tree} vorhergesagt, während die Kartierung dort 15 Prozent Verbuschung nennt. Der größte Teil dieser Baumfläche gehört zum umgebenden Wald und nicht zum Gehölzaufwuchs auf der Offenfläche. Zwei und drei Meter als Grenze schließen umgekehrt das inzwischen herausgewachsene Gehölz aus und unterschätzen den Verbuschungsgrad entsprechend.

Belegen lässt sich das an der Kronenhöhe. Auf der stark verbuschten Fläche 13001\_4443 liegt die mediane Kronenhöhe der als Gehölz vorhergesagten Zellen bei 3,4 Metern, auf der fast offenen Fläche 13001\_2611 dagegen bei 11,6 Metern. Die Grenze bei acht Metern liefert nahezu dieselben Kennzahlen wie die bei fünf, was für eine breite Lücke zwischen aufkommendem Gehölz und Altbestand spricht. Die Wahl der Grenze hängt somit nicht an einem genauen Wert.

\begin{table}[H]
    \centering
    \caption[Kartierter und berechneter Verbuschungsgrad der
    Heidepolygone]{Kartierter und berechneter Verbuschungsgrad der 14 Polygone
    des Lebensraumtyps 4030, sortiert nach dem kartierten Wert. Die Kronenhöhe
    ist der Median der als Gehölz vorhergesagten Zellen.
    Eigene Darstellung.}
    \label{tab:heidepolygone}
    \begin{tabular}{lrrrc}
    \hline
    Polygon & kartiert & berechnet & Kronenhöhe & 2014 gepflegt \\
    \hline
    13001\_4443 & 50 \% & 53,0 \% & 3,4 m  &     \\
    13001\_4444 & 15 \% & 29,2 \% & 2,2 m  &     \\
    13001\_1504 & 15 \% & 1,8 \%  & 1,3 m  &     \\
    13001\_2615 & 10 \% & 2,4 \%  & 5,6 m  &     \\
    13001\_2624 & 2 \%  & 2,3 \%  & 6,4 m  &     \\
    13001\_4442 & 1 \%  & 29,1 \% & 3,8 m  & ja  \\
    13001\_2611 & 1 \%  & 2,6 \%  & 11,6 m &     \\
    13001\_1502 & 1 \%  & 4,6 \%  & 1,9 m  & ja  \\
    13001\_4452 & 0 \%  & 7,3 \%  & 1,9 m  & ja  \\
    13001\_2612 & 0 \%  & 2,5 \%  & 5,5 m  & ja  \\
    13001\_2664 & 0 \%  & 2,8 \%  & 1,1 m  & ja  \\
    13001\_2666 & 0 \%  & 3,6 \%  & 1,0 m  & ja  \\
    13001\_1500 & 0 \%  & 15,1 \% & 1,2 m  & ja  \\
    13001\_4434 & 0 \%  & 13,9 \% & 6,0 m  &     \\
    \hline
    \end{tabular}
\end{table}

Mit dieser Grenze ergeben sich für die Heidepolygone die Werte in Tabelle~\ref{tab:heidepolygone}. Bei den stark verbuschten Flächen stimmen Kartierung und Verfahren gut überein, am deutlichsten auf der Fläche 13001\_4443, deren kartierte 50 Prozent zugleich der höchste Wert sind, den die Kartierung im gesamten Authausener Wald vergibt. Sie gehört zusammen mit 13001\_4444 und 13001\_4442 zur HeideWildnis (Abbildung~\ref{fig:verbuschungskarte_heidewildnis}). Diese drei Flächen tragen mit 50, 15 und einem Prozent die gesamte Spannweite der Referenz und stammen aus einem Gebiet, in dem keine annotierte Teilfläche liegt. Über alle 14 Heidepolygone hinweg beträgt der mittlere Fehler 11,2 Prozentpunkte bei einem Versatz von $+5,4$ Prozentpunkten und einem Bestimmtheitsmaß von 0,28. Der Unterschied zu Tabelle~\ref{tab:hoehengrenze} geht allein auf die sieben Flächen mit Pflegevermerk zurück, die dort nicht enthalten sind.

\begin{figure}[H]
    \centering
    \includegraphics[width=0.75\textwidth]{verbuschungskarte_heidewildnis}
    \caption[Verbuschungsgrad im Gebiet HeideWildnis]{Berechneter
    Verbuschungsgrad im Gebiet HeideWildnis, Zellgröße ein Meter, beschriftet
    mit dem 2014 geschätzten Wert. Eigene Darstellung, Hintergrund: Google
    Satellite.}
    \label{fig:verbuschungskarte_heidewildnis}
\end{figure}

Bei den schwach verbuschten Flächen geht diese Übereinstimmung verloren. Bei den sieben der 14 Flächen, die einen Pflegevermerk tragen, liegt der berechnete Wert ausnahmslos über dem kartierten. Auch ohne diese Flächen bleibt die Rangkorrelation nach Spearman mit 0,23 niedrig, das Verfahren bringt die schwach verbuschten Flächen also nicht in die richtige Reihenfolge. Die Ursache wird in Abschnitt~\ref{diskussion_teil} behandelt.

Unabhängig von den Referenzflächen liefert das Verfahren für jede Rasterzelle des Untersuchungsgebiets einen Wert. Über alle vier Befliegungen sind das 64,5 Hektar oder rund 25.800 Zellen. Die Hälfte dieser Zellen liegt unter 4,2 Prozent Verbuschung, der Mittelwert beträgt dagegen 20,0 Prozent, denn gut ein Viertel der Zellen überschreitet 30 Prozent und zieht ihn nach oben. Wie Abbildung~\ref{fig:verbuschungskarte_uebersicht} zeigt, stehen große offene Bereiche einzelnen, stark verbuschten Teilflächen gegenüber. Eine einzelne Kennzahl beschreibt eine so verteilte Fläche deshalb schlecht.

\begin{figure}[H]
    \centering
    \includegraphics[width=\textwidth]{verbuschungskarte_uebersicht}
    \caption[Verbuschungsgrad der vier Befliegungsgebiete]{Berechneter
    Verbuschungsgrad der vier Befliegungsgebiete, Zellgröße 20 Meter. Eigene
    Darstellung, Hintergrund: Google Satellite.}
    \label{fig:verbuschungskarte_uebersicht}
\end{figure}


\subsection{Diskussion}\label{diskussion_teil}

Die bisherigen Abschnitte haben die Zahlen berichtet, die folgenden ordnen sie ein. Im Mittelpunkt steht dabei die Frage, welcher Anteil der Abweichungen dem Modell zuzuschreiben ist und welcher der Referenz oder dem Verfahren, aus dem die Labels stammen.

\subsubsection{Bedeutung der Höheninformation}

Den größten Einzeleffekt in der gesamten Auswertung hat nicht die Wahl des
Modells, sondern die Höhengrenze. Sie senkt den mittleren Fehler von 16,9 auf
9,5 Prozentpunkte und hebt das Bestimmtheitsmaß von einem negativen Wert auf
0,65. Der Grund liegt darin, dass die Klasse \texttt{tree} Pioniergehölz und
Altbestand nicht trennt (Abschnitt~\ref{verbuschungsmass}). Bestätigt wird das durch die Kartierung von 2014, die
13001\_4443 als \enquote{Stark von Birke (50 \%, bis 3 m hoch) verbuschte
Callunaheide} beschreibt, während die photogrammetrisch gemessene Kronenhöhe
dort im Median bei 3,4 Metern liegt. Aus der Klasse allein ist diese
Unterscheidung nicht zu gewinnen, da beide Bestände dieselbe Klasse tragen. Erst die
Höhe je Zelle macht aus der Segmentierung eine Größe, die der fachlichen
Definition entspricht, nach der Verbuschung das Aufkommen von Pioniergehölzen
auf offener Fläche meint und nicht den benachbarten Wald.

\subsubsection{Zeitlicher Abstand zwischen Kartierung und Befliegung}

Zwischen Kartierung und Befliegung liegen elf Jahre. Verbuschung ist genau
der Prozess, der sich in dieser Zeit verändert. Sieben der 14 Heidepolygone
tragen einen Vermerk über eine kurz zuvor durchgeführte Gehölzentnahme. Gerade
bei ihnen liegt der berechnete Wert ausnahmslos über dem kartierten, am
deutlichsten auf der Fläche 13001\_1500, für die die Kartierung eine kürzlich
erfolgte Freistellung unter der Stromleitung vermerkt und das Verfahren eine
Verbuschung von 15,1 Prozent ausweist.

Der im Feld \texttt{VERBUP} erhobene Wert von null Prozent war im Jahr 2014
durchaus richtig, für 2025 gilt er jedoch nicht mehr. An den Stellen, wo das Verfahren mehr Verbuschung ausweist als die Kartierung, ist deshalb nicht zwangsläufig das Modell im Unrecht, sondern in mehreren Fällen die Referenz veraltet. Aus demselben Grund bleibt auch die Wahl der Höhengrenze ein Kompromiss, weil in elf Jahren herangewachsenes Gehölz genau in dem Bereich liegt, den die Grenze abschneidet.

\subsubsection{Vergleich mit der Geländekartierung}

Ob ein mittlerer Fehler von 9,5 Prozentpunkten gut ist, zeigt erst der
Vergleich mit publizierten Arbeiten. Schmidt et al. erreichen bei der
fernerkundlichen Schätzung der Calluna-Deckung ein Bestimmtheitsmaß von 0,79
bei einem mittleren Fehler von 15,7 Prozentpunkten, gestützt auf 700
Stichprobenpunkte \cite{Schmidt.2017}. Die hier erreichten Werte liegen in derselben Größenordnung, beruhen jedoch auf sieben Flächen. Der dortige Wertebereich reicht mit null bis 96 Prozent etwa doppelt so weit. Aussagekräftiger als der Mittelwert ist ohnehin die Verteilung der Abweichungen, denn drei der sieben Flächen liegen innerhalb von drei Prozentpunkten, drei weitere weichen um 13 bis 14 Punkte ab. Der Versatz von $+1,8$ Prozentpunkten zeigt zugleich, dass sich diese Abweichungen weitgehend ausgleichen und das Verfahren nicht systematisch über- oder
unterschätzt.

Wie stark die Übereinstimmung bei den kräftig verbuschten Flächen von einer einzigen
getragen wird, zeigt eine Kontrollrechnung ohne 13001\_4443, bei der die
Korrelation von 0,86 auf 0,32 fällt. Bei sieben Beobachtungen und einem Wert,
der mehr als dreimal so hoch liegt wie alle übrigen, ist das zu erwarten. Dass die gegen Ausreißer robuste Rangkorrelation mit 0,23 weit unter dem Bestimmtheitsmaß von 0,65 liegt, zeigt, dass die Übereinstimmung an den Absolutwerten einzelner Flächen hängt und nicht an einer durchgängig richtigen Rangfolge.

Erklären lassen sich die größten Abweichungen über die Bemerkungsfelder. Für
die Fläche 13001\_2615 nennt die Kartierung eine beginnende Verbuschung mit
Birke und Traubenkirsche, also kleine, verstreute Einzelgehölze mit winziger
Kronenprojektion, die ein Kartierer im Gelände dennoch als deutliches Signal
wahrnimmt. Für 13001\_4434 beschreibt sie ein Übergangsbiotop zwischen
verheidetem Borstgrasrasen und beerstrauchreichem Kiefernwald, wo der
Gehölzanteil als Wesenszug des Biotoptyps gilt und nicht als Verbuschung. Darin
zeigt sich, dass Referenz und Verfahren bei den schwach verbuschten Flächen verschiedene Dinge messen, denn der Kartierer bewertet, ob der Prozess eingesetzt hat, wohingegen das Verfahren überschirmte Fläche misst. Das ist keine Schwäche des
Modells, sondern eine inhaltliche Grenze des Vergleichs.

\subsubsection{Segmentierungsgüte und Herkunft der Labels}

Einordnen lässt sich die erreichte Güte über veröffentlichte Ergebnisse
derselben Architektur. Auf dem Straßenraumdatensatz Toronto3D erreicht
PointNet++ eine mittlere IoU von 0,418 \cite{Diab.2022} und liegt damit dicht
bei den hier erzielten 0,393, wobei sich beide Datensätze in Sensorik und
Klassenschema unterscheiden und deshalb nur die Größenordnung vergleichbar ist.
Der Mittelwert verdeckt dabei, wie ungleich sich die Güte über die Klassen
verteilt, denn großflächige und strukturell eindeutige Klassen wie Straße und
Vegetation erreichen dort 0,893 und 0,690, wohingegen Fahrbahnmarkierungen und
Zäune bei 0,000 und 0,030 liegen. Die schwachen Werte für seltene Klassen wie
\texttt{deadwood} und \texttt{ground} folgen demnach einem bekannten Muster
dieser Architektur.

Anders verhält es sich mit der Verwechslung zwischen Zwergstrauchheide, Gras
und Busch, die allen fünf Ansätzen der Ablationsstudie standgehalten hat und
erst auf die Nachbarschaftsradien reagierte (Abschnitt~\ref{ergebnis_ablation}).

Neben der Punktdichte kommt die Herkunft der Labels als Erklärung in Betracht,
denn der Verschnitt überträgt die Klasse eines Polygons auf die gesamte
Punktsäule darunter (Abschnitt~\ref{label_transfer}). Diesen Fehler in beide Richtungen greift die höhenbasierte Korrektur auf, die sich dabei allerdings eine eigene Schwäche einhandelt, weil sie punktweise und nicht objektbezogen arbeitet. Ein Strauch beispielsweise, dessen Basis am Boden ansetzt, wird an der Grenze von 0,5 Metern waagerecht durchgeschnitten. Hinzu kommt die Reihenfolge der Regeln, denn 40,6 Prozent aller Punkte, die neu die Klasse \texttt{bush} erhalten, durchlaufen zuvor eine Hochstufung zu \texttt{tree} und werden erst von der Abwärtsregel wieder zurückgestuft.

Daraus folgt eine These, die sich mit den vorliegenden Daten nicht abschließend
prüfen lässt, als Erklärungsansatz aber gut begründet ist. Das Netz wird auf
genau diesen Labels trainiert und lernt die Grenze bei 0,5 Metern, als wäre sie
eine physische Eigenschaft der Vegetation. Ein Teil der Verwechslung zwischen
\texttt{heath} und \texttt{bush}, die sich als hartnäckigste Schwäche des
Modells zeigt, wäre dann ein Artefakt der Labelvergabe und keine Grenze der
Architektur.

\subsubsection{Limitationen}

Mit sieben Heidepolygonen in der bereinigten Auswahl und 14 einschließlich der
gepflegten Flächen ist die Stichprobe zu klein, um Bestimmtheitsmaß und
Korrelation stabil zu schätzen, was der Einbruch der Korrelation nach dem
Entfernen einer einzigen Fläche unmittelbar zeigt. Hinzu kommt, dass die
Referenz selbst auf grob gestuften Augenmaßschätzungen im Gelände beruht.

Schwerer wiegt ein Ungleichgewicht zwischen Bewertung und Anwendung, denn
Validierungs- und Testteil bestehen zu 83 beziehungsweise 66 Prozent aus
Kacheln der Cuxhavener Küstenheide, während der Verbuschungsgrad ausschließlich
für die Authausener Flächen berechnet wurde. Bewertet wird das Modell also
überwiegend auf einem Landschaftstyp, für den keine Verbuschungsberechnung
vorliegt. Der Validierungsteil kann zudem die letzten drei Modellvarianten
nicht unterscheiden, sodass die Modellauswahl zu einem Teil zufällig ausfällt
(Abschnitt~\ref{ergebnis_ablation}).

Schließlich beruhen die Schwellenwerte der Labelkorrektur auf den 16
annotierten Teilflächen und wären bei einer Erweiterung des Datensatzes neu zu
bestimmen. In dicht bewaldeten Bereichen rekonstruiert Structure from Motion
außerdem fast nur die Kronenoberfläche, sodass die Höhennormalisierung dort
unsicher wird. Für die Offenlandbiotope dieser Arbeit ist das unkritisch, für
eine flächendeckende Karte des gesamten Waldgebietes wäre es zu
berücksichtigen.

\subsubsection{Anwendbarkeit in der Praxis}

Als Ausgangsproblem nennt die Einleitung, dass eine Geländekartierung zeit- und
kostenintensiv ist und deshalb immer nur einen Bruchteil der Fläche abdeckt.
Genau dieses Verhältnis kehrt das Verfahren um, denn es liefert Werte für 64,5
Hektar, von denen sich nur 9,9 überprüfen lassen, und macht den
Verbuschungsgrad für jede beliebige Bezugsfläche berechenbar. Hinzu kommt die
Wiederholbarkeit, denn eine erneute Befliegung liefert mit demselben Modell
unmittelbar eine Veränderungsanalyse, was eine einmalige Geländekartierung
nicht leistet, für die Erfolgskontrolle von Pflegemaßnahmen aber erforderlich
wäre.

Ob Naturschutzpraktiker das Verfahren selbst anwenden können, war nicht mehr
Gegenstand dieser Arbeit, die Voraussetzungen dafür lassen sich jedoch
benennen. Sämtliche Ergebnisse liegen in gängigen, georeferenzierten Formaten
vor, die klassifizierten Punktwolken als LAZ und die Karten als GeoTIFF in
EPSG~25832. Da QGIS seit Version 3.18 Punktwolken nativ darstellt und sich
dafür auf die Bibliothek PDAL stützt \cite{QGIS.318}, lassen sich die
Ergebnisse ohne Zusatzsoftware betrachten und weiterverwenden.

Offen bleibt allein die Inferenz innerhalb der GIS-Umgebung. Für
zweidimensionale Netze steht dafür mit Deepness ein etabliertes QGIS-Plugin
bereit, das Modelle im Format ONNX (Open Neural Network Exchange) ausführt
\cite{Aszkowski.2023}. Es arbeitet jedoch ausschließlich auf Rasterlayern und
ist auf ein punktbasiertes Netz wie PointNet++ nicht anwendbar. Da ein
Gegenstück für Punktwolken bislang fehlt, laufen Vorverarbeitung und Inferenz
vorerst außerhalb von QGIS und nur die Ergebnisse werden dort eingebunden.
Weil die Verarbeitungsschritte bereits skriptgesteuert und über
Umgebungsvariablen parametrisiert sind, würde eine spätere Einbindung
keine methodische Änderung erfordern, sondern lediglich eine Anbindung an die
Verarbeitungsumgebung von QGIS.

Für die Praxis folgt daraus eine Aufgabenteilung. Stark verbuschte Flächen
trennt das Verfahren zuverlässig von offenen, sodass sich damit eine Rangfolge
für anstehende Pflegemaßnahmen aufstellen lässt. Schwach verbuschte Flächen, etwa solche zwischen einem und zehn Prozent,
stuft es dagegen nicht zuverlässig ab. Die Geländekartierung ersetzt es dort nicht, wohl
aber weist es ihr die Flächen zu, auf denen sie sich lohnt.
